\section*{Chapter \ref{ch:pl:time-series-databases-and-the-alternatives} --- Time-Series Databases and the Alternatives}

\begin{solution}{exo:pl:time-series-databases-and-the-alternatives:1}
Point lookup: $3.8 / 0.05$, about 80 times faster. Cross-section: $313 / 7$, over 40 times slower. (The exact factors move from one
measurement to the next; their size does not.)
\end{solution}

\begin{solution}{exo:pl:time-series-databases-and-the-alternatives:2}
The index is a B-tree ordered by (symbol, time): the lookup descends a few pages to the symbol's last entry before the time and
reads one row. The cross-section needs every row of every symbol; the index offers no shortcut, and the row store reads each whole
row one at a time to use two of its fields.
\end{solution}

\begin{solution}{exo:pl:time-series-databases-and-the-alternatives:3}
$6.5 \times 60 = 390$ buckets. S007 had no quote in 31 of them: \omterm{def:pl:time-series-databases-and-the-alternatives:bucket}{time buckets} without data produce no row, and a bar series that
needs every minute must fill the empty ones explicitly (carrying the last mid forward, with a zero count).
\end{solution}

\begin{solution}{exo:pl:time-series-databases-and-the-alternatives:4}
A machine's speed drifts during a benchmark (frequency, other jobs, the page cache). Running one engine seven times and then the
next would attribute the drift to the engines; interleaving exposes every engine to the same conditions in every round.
\end{solution}

\begin{solution}{exo:pl:time-series-databases-and-the-alternatives:5}
Both are dominated by the fixed cost $c_0$ (opening the day's file, reading its footer, planning). And the point query, as written
(the last quote at or before noon, ordered by time, limit one), touches every quote of the symbol up to noon, about half its day,
and sorts them, while the range touches only the hour's 153 rows, pruned by \omterm{def:pl:columnar-formats:pushdown}{zone maps}. Rows touched, not rows returned, set the
variable cost.
\end{solution}

\begin{solution}{exo:pl:time-series-databases-and-the-alternatives:6}
Equal cost when $(1-f)\,3.8 + 313 f = (1-f)\,0.05 + 7 f$ (in the measured medians' units): $f = (3.8-0.05)/[(3.8-0.05) +
(313-7)] = 1.2\%$.
\end{solution}

\begin{solution}{exo:pl:time-series-databases-and-the-alternatives:7}
They agree on all five shapes: one row for the point lookup, 171 quotes in the hour, 358 one-minute bars, 2\,000 trades joined and
50 average spreads.
\end{solution}

\begin{solution}{exo:pl:time-series-databases-and-the-alternatives:8}
One query shape, the vendor's, may not be the firm's workload; one run each, in sequence, confounds the engines with the page cache
(the second run finds the data warm) and with the machine's drift; nothing checked that both gave the same answer. Benchmark the
firm's own shapes on its own data, with answers checked, interleaved rounds, and cold and warm reported apart.
\end{solution}

\begin{solution}{pb:pl:time-series-databases-and-the-alternatives:1}
\begin{enumerate}
\item Point (the last quote before an order), range (one symbol's hour), bars (minute bars of one day), as-of join (each trade
with its quote), cross-section (average spread by symbol over a month).
\item Point, range and bars touch few rows; the as-of join touches two slices; the cross-section touches all.
\item Sorted by symbol then time, with \omterm{def:pl:columnar-formats:pushdown}{zone maps} (or an index), so that one symbol's window is a contiguous slice.
\item Packaging of the same ideas: a server that ingests, compacts, retains and answers, with time-specific operations built in.
\item Replacing a bucket's records by a summary; in the \omterm{def:pl:capturing-and-storing-tick-data:tier}{retention policy}, ticks age out after months while bars stay for years.
\item An embedded engine runs inside the querying process (no round trip, no serialisation) but gives up concurrent writers,
access control and shared service; a server gives those and pays the round trip.
\item Running each operator on a batch of a column's values, so call overhead is paid per batch and inner loops are tight array
loops.
\item Carrying positions of qualifying rows through the filters and reading the other columns only for the rows that survive.
\item A B-tree walks a few pages to one row; a scan of a file still pays opening, metadata and decoding a block.
\item Fastest: SQLite on the point lookup, numpy on the range, the bars and the cross-section, numpy and SQLite level on the as-of
join; slowest: pandas on four shapes, SQLite on the cross-section.
\item 1, 153, 359, 2\,000 and 50; because a fast wrong answer is the commonest benchmark result, and five engines disagreeing would
make the timings meaningless.
\item SQLite 0.05~ms and 313~ms, DuckDB 3.8~ms and 7~ms.
\item It is written for this layout and these queries (sorted arrays, an offset per symbol); every new question needs new code and
new tests, which a general engine provides for free.
\item The fresh engine's start-up costs and whatever was not yet cached, which the first query of a session pays.
\item The machine, its load, the page-cache state, the data size and layout, and the engines' versions.
\item \textbf{Named result.} On the month of quotes, SQLite's index wins the point lookup (0.05~ms against DuckDB's 3.8), the
columnar engines win the scans (the cross-section in 7~ms against SQLite's 313), a hand-written engine beats both where it applies,
and DuckDB beats SQLite on a mixed workload as soon as more than about 1.2\% of the queries are cross-sections.
\item DuckDB (or Polars): the workload is scans and aggregations over history.
\item An in-process index (SQLite, or a sorted in-memory array with offsets, as in the numpy engine), next to the router: a point
lookup must cost microseconds, not a file open and a plan.
\item When many writers and readers share live data, need access control and retention without the firm building it, and the
benchmark on the firm's shapes shows it wins.
\item The shape of the query: how many rows it must touch against the engine's fixed and per-row costs.
\end{enumerate}
\end{solution}

\begin{solution}{iq:pl:time-series-databases-and-the-alternatives:1}
A database optimised for time-keyed records appended in time order and read by time ranges, with \omterm{def:pl:time-series-databases-and-the-alternatives:bucket}{time buckets}, as-of joins and
retention built in. For research on \omterm{def:pl:capturing-and-storing-tick-data:tick}{tick data}, files in open columnar formats and an embedded engine often suffice; a server is
worth it when many writers and readers share live data.
\iqlookfor{The category is packaging of known ideas; the query mix decides.}
\end{solution}

\begin{solution}{iq:pl:time-series-databases-and-the-alternatives:2}
It reads only the needed columns, in batches, with tight vectorised loops and \omterm{def:pl:time-series-databases-and-the-alternatives:vectorised}{late materialisation}, so scans are cheap per row.
A single row still costs a fixed amount (open, plan, decode a block), which an index on a row store avoids.
\iqlookfor{Fixed against per-row cost.}
\end{solution}

\begin{solution}{iq:pl:time-series-databases-and-the-alternatives:3}
Use your own data and query shapes; check both return identical answers first; interleave runs; report cold and warm separately
with spreads; state machine, load and cache; vary the data size to see how each scales.
\iqlookfor{Correctness before speed; interleaving; cold and warm.}
\end{solution}

\begin{solution}{iq:pl:time-series-databases-and-the-alternatives:4}
When the data are files one team reads (a \omterm{def:pl:a-tick-store-on-open-formats:store}{tick store}, research datasets) and queries are analytical: no server to run, no
round trips, zero-copy into dataframes. A server is needed for many concurrent writers, access control, or one shared live
dataset.
\iqlookfor{Who writes, and how many readers share it.}
\end{solution}

\begin{solution}{iq:pl:time-series-databases-and-the-alternatives:5}
Inventory the queries it answers (shapes, frequencies, latency needs) and its writers; build candidates on open formats (files
and embedded engines, and one or two servers); benchmark them on the firm's data with answers checked against the old system;
migrate readers shape by shape, keeping both until the answers match for weeks.
\iqlookfor{Workload inventory, a checked benchmark, parallel running.}
\end{solution}
